\section{Wick's theorem and the Gaussian class}
\label{sec:wick}
\motiv{A state is \emph{Gaussian} when all its Majorana moments are given by Wick's rule
from its covariance. The rotated product mixtures $U^*\rho_\lambda U$ of
Section~\ref{sec:hamiltonian} are Gaussian (Theorem~\ref{thm:mixture-gaussian}), and with
them the Gibbs and thermal states \vet{and the pure states $U^*|s\rangle\langle s|U$}. Since the Majorana monomials form a basis of the
operators on $\mathcal F_n$, a Gaussian state is determined by its covariance
(Theorem~\ref{thm:unique}); hence every Gaussian state is a rotated product mixture, and the
Gibbs state with a given covariance is unique.}

The conventions of Sections~\ref{sec:hamiltonian} and~\ref{sec:thermal} are in force.
Vectors $v\in\C^{2n}$ are indexed by the labels $x_1,\ldots,x_n,y_1,\ldots,y_n$; $e_a$ is the
unit vector of the label $a$, and $v\cdot c=\sum_av_ac_a$. The annihilation operator of mode
$j$ is $a_j=\tfrac12(c_{x_j}+\ii c_{y_j})$, so that $a_j=u\cdot c$ with
$u=\tfrac12(e_{x_j}+\ii e_{y_j})$. For a density operator $\rho$,
$\tr(\rho c_ac_b)=\delta_{ab}-\ii\Gamma(\rho)_{ab}$ (Theorem~\ref{thm:reality} and
$c_ac_b=\delta_{ab}+\tfrac12[c_a,c_b]$), which motivates the \emph{two-point matrix}
\[
 \G(\rho)=1-\ii\,\Gamma(\rho),
\]
a complex matrix on the labels with $\G(\rho)_{ab}=\tr(\rho c_ac_b)$.

\subsection{Wick's rule and the Gaussian class}
\label{sec:gaussian-class}

\begin{definition}[Wick expansion, Gaussian state]\label{def:wick}
For a complex matrix $\G$ on the labels and vectors $v_1,\ldots,v_m\in\C^{2n}$, the
\emph{Wick expansion} is defined by $\W_\G()=1$ and, for $m\geq1$,
\[
 \W_\G(v_1,\ldots,v_m)=\sum_{j=2}^m(-1)^j\,(v_1^{\T}\G v_j)\,
 \W_\G(v_2,\ldots,\widehat{v_j},\ldots,v_m),
\]
where the hat removes the entry. A density operator $\rho$ is \emph{Gaussian} when
\[
 \tr\bigl(\rho\,(v_1\cdot c)(v_2\cdot c)\cdots(v_m\cdot c)\bigr)=\W_{\G(\rho)}(v_1,\ldots,v_m)
 \qquad\text{for all }m\geq0\text{ and all }v_1,\ldots,v_m\in\C^{2n}.
\]
\end{definition}
For $m=1$ the sum is empty and $\W_\G(v_1)=0$; by induction $\W_\G$ vanishes for every odd
$m$. For $m=2$, $\W_\G(v_1,v_2)=v_1^{\T}\G v_2$, and for $m=4$,
\[
 \W_\G(v_1,v_2,v_3,v_4)=(v_1^{\T}\G v_2)(v_3^{\T}\G v_4)-(v_1^{\T}\G v_3)(v_2^{\T}\G v_4)
 +(v_1^{\T}\G v_4)(v_2^{\T}\G v_3),
\]
the sum over pairings with the sign of the pairing. \revf{For every even $m$, $\W_\G(v_1,\ldots,v_m)$ is the Pfaffian
of the antisymmetric $m\times m$ matrix whose entries above the diagonal are
$v_i^{\T}\G v_j$~\cite{bravyi}: the recursion of Definition~\ref{def:wick} is the expansion of
that Pfaffian along its first row.} With $v_i=e_{a_i}$ the rule reads
$\tr(\rho\,c_{a_1}\cdots c_{a_m})=\W_{\G(\rho)}(e_{a_1},\ldots,e_{a_m})$, Wick's rule in its usual
form with the pairing $\tr(\rho c_ac_b)=\delta_{ab}-\ii\Gamma_{ab}$; repeated labels are allowed.
We state the rule for vectors rather than labels because a Gaussian unitary acts on the
vectors, $U(v\cdot c)U^*=(O^{\T}v)\cdot c$ (Theorem~\ref{thm:transport}), and because the
relation that drives Theorem~\ref{thm:mixture-gaussian} involves the combinations
$c_{x_j}\pm\ii\lambda_jc_{y_j}$; the label form is the special case $v_i=e_{a_i}$.

\begin{lemma}[Multilinearity and the two-point rule]\label{lem:wick-basic}
\begin{enumerate}
\item[(a)] $\W_\G(v_1,\ldots,v_m)=0$ for odd $m$, $\W_\G$ is linear in $v_1$, and
$\W_\G(v,w)=v^{\T}\G w$.
\item[(b)] For every density operator $\rho$ and all $v,w$,
$\tr\bigl(\rho\,(v\cdot c)(w\cdot c)\bigr)=v^{\T}\G(\rho)w$; in particular Wick's rule holds
for $m=0$ and $m=2$ for every density operator.
\item[(c)] $(u\cdot c)(w\cdot c)+(w\cdot c)(u\cdot c)=2(u^{\T}w)\,1$ for all $u,w\in\C^{2n}$.
\end{enumerate}
\end{lemma}
\begin{proof}
(a) We noted odd vanishing above; linearity in $v_1$ is linearity of $v_1^{\T}\G v_j$ in each
term; $m=2$ is the definition with one term, $j=2$.
(b) $\tr(\rho(v\cdot c)(w\cdot c))=\sum_{ab}v_aw_b\tr(\rho c_ac_b)=\sum_{ab}v_a\G(\rho)_{ab}w_b$;
for $m=0$ both sides are $\tr\rho=1$.
(c) By Theorem~\ref{thm:car},
$\sum_{ab}u_aw_b(c_ac_b+c_bc_a)=\sum_{ab}u_aw_b\,2\delta_{ab}=2(u^{\T}w)$.
\end{proof}
\begin{remark}\label{rem:wick-condition}
Beyond $m=2$, Wick's rule is a condition. For $n=1$, $\rho=\tfrac12(1+c_{x_1})$ is a
density operator ($c_{x_1}$ is Hermitian with $c_{x_1}^2=1$) with $\tr(\rho c_{x_1})=1\neq0$. For
$n=2$, $M=c_{x_1}c_{y_1}c_{x_2}c_{y_2}$ is Hermitian with $M^2=1$, so
$\rho=\tfrac14(1+M)$ is a density operator; $\Gamma(\rho)=0$ since $c_ac_bM$ is $\pm$ a monomial
of two labels for $a\neq b$ (Lemma~\ref{lem:monomial-trace}), while
$\tr(\rho M)=1\neq0=\W_1(e_{x_1},e_{y_1},e_{x_2},e_{y_2})$. Theorem~\ref{thm:mixture-gaussian}
proves the rule for the product mixtures.
\end{remark}

\begin{theorem}[Transport by a Gaussian unitary]\label{thm:transport}
Let the unitary $U$ implement $O$ and let $\rho$ be a density operator. Then
$U(v\cdot c)U^*=(O^{\T}v)\cdot c$ for every $v$,
$\W_\G(O^{\T}v_1,\ldots,O^{\T}v_m)=\W_{O\G O^{\T}}(v_1,\ldots,v_m)$ for every $\G$, and
$\G(U^*\rho U)=O\,\G(\rho)\,O^{\T}$. If $\rho$ is Gaussian, so is $U^*\rho U$.
\end{theorem}
\begin{proof}
$U(v\cdot c)U^*=\sum_av_a\sum_bO_{ab}c_b=\sum_b(O^{\T}v)_bc_b$. The identity for $\W$ follows
by induction on $m$ from $(O^{\T}v_1)^{\T}\G(O^{\T}v_j)=v_1^{\T}(O\G O^{\T})v_j$. The two-point
matrix transforms with the covariance, $\Gamma(U^*\rho U)=O\Gamma(\rho)O^{\T}$
(Theorem~\ref{thm:congruence}), and $O\,1\,O^{\T}=1$. Finally, with $X=\prod_i(v_i\cdot c)$,
cyclicity of the trace gives
\begin{align*}
 \tr(U^*\rho U\,X)&=\tr(\rho\,UXU^*)=\tr\Bigl(\rho\prod_i\bigl((O^{\T}v_i)\cdot c\bigr)\Bigr)\\
 &=\W_{\G(\rho)}(O^{\T}v_1,\ldots,O^{\T}v_m)=\W_{O\G(\rho)O^{\T}}(v_1,\ldots,v_m),
\end{align*}
which is $\W_{\G(U^*\rho U)}(v_1,\ldots,v_m)$.
\end{proof}

Removing one factor at a time reduces Wick's rule to two identities
(Lemma~\ref{lem:move}) and a criterion (Lemma~\ref{lem:criterion}).

\begin{lemma}[Moving a factor through a product]\label{lem:move}
Let $b=u\cdot c$, $X=(w_1\cdot c)\cdots(w_k\cdot c)$, and let $X_{\hat\imath}$ be the product
with the $i$-th factor removed. Then
\[
 bX=\sum_{i=1}^k(-1)^{i-1}\,2(u^{\T}w_i)\,X_{\hat\imath}+(-1)^kXb,\qquad
 Xb=\sum_{i=1}^k(-1)^{k-i}\,2(u^{\T}w_i)\,X_{\hat\imath}+(-1)^kbX.
\]
\end{lemma}
\begin{proof}
Induction on $k$; for $k=0$ both identities read $b=b$. Write $X=(w_1\cdot c)X'$ with
$X'=(w_2\cdot c)\cdots(w_k\cdot c)$, so that $X_{\hat1}=X'$ and
$X_{\hat\imath}=(w_1\cdot c)X'_{\widehat{i-1}}$ for $i\geq2$. By Lemma~\ref{lem:wick-basic}(c),
$b(w_1\cdot c)=2(u^{\T}w_1)-(w_1\cdot c)b$, so $bX=2(u^{\T}w_1)X'-(w_1\cdot c)\,bX'$; the
induction hypothesis $bX'=\sum_{i=2}^k(-1)^{i-2}2(u^{\T}w_i)X'_{\widehat{i-1}}+(-1)^{k-1}X'b$,
multiplied by $-(w_1\cdot c)$, turns the signs into $(-1)^{i-1}$ and $(-1)^k$ and the products
into $X_{\hat\imath}$ and $Xb$: the first identity. The second is the first solved for $Xb$:
multiply the first identity by $(-1)^k$ and use $(-1)^k(-1)^{i-1}=-(-1)^{k-i}$.
\end{proof}

\begin{lemma}[A recursion criterion]\label{lem:criterion}
Let $\rho$ be a density operator such that (i) $\tr(\rho X)=0$ for every product $X$ of an odd
number of factors $w\cdot c$, and (ii) for every label $a$ and every product
$X=(w_1\cdot c)\cdots(w_k\cdot c)$ with $k$ odd,
\[
 \tr(\rho\,c_aX)=\sum_{i=1}^k(-1)^{i-1}\,\bigl(e_a^{\T}\G(\rho)w_i\bigr)\,\tr(\rho X_{\hat\imath}).
\]
Then $\rho$ is Gaussian.
\end{lemma}
\begin{proof}
We prove Wick's rule by induction on $m$. For $m$ odd both sides vanish, by (i) and
Lemma~\ref{lem:wick-basic}(a). For $m$ even and $m\geq2$, write $v_1\cdot c=\sum_av_{1a}c_a$
and apply (ii) with $X=(v_2\cdot c)\cdots(v_m\cdot c)$, $k=m-1$ odd:
\[
 \tr\bigl(\rho\,(v_1\cdot c)X\bigr)=\sum_av_{1a}\sum_{i=1}^{m-1}(-1)^{i-1}(e_a^{\T}\G w_i)
 \tr(\rho X_{\hat\imath})=\sum_{j=2}^m(-1)^j(v_1^{\T}\G v_j)\,\tr(\rho X_{\widehat{j-1}}),
\]
with $w_i=v_{i+1}$, $j=i+1$, $(-1)^{i-1}=(-1)^j$. By the induction hypothesis
$\tr(\rho X_{\widehat{j-1}})=\W_\G(v_2,\ldots,\widehat{v_j},\ldots,v_m)$, and the right side is
$\W_\G(v_1,\ldots,v_m)$ by Definition~\ref{def:wick}. For $m=0$ both sides are $1$.
\end{proof}

\subsection{The product mixtures are Gaussian}
\label{sec:mixture-gaussian}

\begin{lemma}[Parity]\label{lem:parity}
Let $\Pi_n=Z_1Z_2\cdots Z_n=\operatorname{diag}\bigl(\prod_j(-1)^{s_j}\bigr)$. Then $\Pi_n^2=1$,
$\Pi_nc_a\Pi_n=-c_a$ for every label $a$, and $\Pi_n\rho_\lambda\Pi_n=\rho_\lambda$ for every
$\lambda\in[-1,1]^n$. Consequently $\tr(\rho_\lambda X)=0$ for every product $X$ of an odd number
of factors $w\cdot c$.
\end{lemma}
\begin{proof}
$\Pi_n$ is diagonal with entries $\pm1$, so $\Pi_n^2=1$ and $\Pi_n$ commutes with the diagonal
matrix $\rho_\lambda$. Writing $c_a=\operatorname{diag}(f_a)X_j$ with $j$ the mode of $a$
(equation~\eqref{eq:signed-flip}) and $\pi(s)=\prod_k(-1)^{s_k}$ for the diagonal entry of $\Pi_n$,
$\Pi_nc_a\Pi_n=\operatorname{diag}(\pi\,f_a)X_j\operatorname{diag}(\pi)=\operatorname{diag}(\pi\,f_a\,(\pi\circ\mathrm{flip}_j))X_j$
and $\pi(\mathrm{flip}_js)=-\pi(s)$, so $\Pi_nc_a\Pi_n=-c_a$. For $X$ of odd length,
$\Pi_nX\Pi_n=(-1)^{\text{length}}X=-X$, by induction on the number of factors: $\Pi_n(w\cdot c)X'\Pi_n
=(\Pi_n(w\cdot c)\Pi_n)(\Pi_nX'\Pi_n)=-(w\cdot c)\,\Pi_nX'\Pi_n$. Hence
$\tr(\rho_\lambda X)=\tr(\Pi_n\rho_\lambda\Pi_n\,X)=\tr(\rho_\lambda\,\Pi_nX\Pi_n)=-\tr(\rho_\lambda X)$.
\end{proof}

\begin{lemma}[The mode relation]\label{lem:mode}
Let $\lambda\in[-1,1]^n$ and fix a mode $j$. Then
\[
 (c_{x_j}-\ii\lambda_jc_{y_j})\,\rho_\lambda=\rho_\lambda\,(c_{x_j}+\ii\lambda_jc_{y_j}),\qquad
 (c_{y_j}+\ii\lambda_jc_{x_j})\,\rho_\lambda=\rho_\lambda\,(c_{y_j}-\ii\lambda_jc_{x_j}).
\]
Moreover $\G(\rho_\lambda)=1-\ii\D(\lambda)$ has the rows
$(\G w)_{x_j}=w_{x_j}-\ii\lambda_jw_{y_j}$ and $(\G w)_{y_j}=w_{y_j}+\ii\lambda_jw_{x_j}$ for every
vector $w$.
\end{lemma}
\begin{proof}
$\rho_\lambda=\operatorname{diag}(p)$ with $p(s)=\prod_kp_k(s_k)$ and
$p_j(s_j)=\tfrac12(1-\lambda_j(-1)^{s_j})$ (Definition~\ref{def:mixture}). The factors
$p_k(s_k)$, $k\neq j$, are common to $p(s)$ and $p(\mathrm{flip}_js)$, so
\begin{equation}\label{eq:weight-flip}
 p_j(s_j)\,p(\mathrm{flip}_js)=p_j(1-s_j)\,p(s)\qquad\text{for every }s,
\end{equation}
that is, $(1-\lambda_jZ_j)\operatorname{diag}(p\circ\mathrm{flip}_j)=\rho_\lambda(1+\lambda_jZ_j)$,
since $2p_j(s_j)=1-\lambda_j(-1)^{s_j}$ is the entry of $1-\lambda_jZ_j$ and $2p_j(1-s_j)$ that of
$1+\lambda_jZ_j$. As $Z_{<j}$ is diagonal and $X_j\operatorname{diag}(p)=\operatorname{diag}(p\circ\mathrm{flip}_j)X_j$
(Lemma~\ref{lem:pauli}),
\[
 (1-\lambda_jZ_j)\,c_{x_j}\rho_\lambda=(1-\lambda_jZ_j)Z_{<j}\operatorname{diag}(p\circ\mathrm{flip}_j)X_j
 =\rho_\lambda(1+\lambda_jZ_j)Z_{<j}X_j=\rho_\lambda(1+\lambda_jZ_j)\,c_{x_j}.
\]
Since $Z_jc_{x_j}=-c_{x_j}Z_j=\ii c_{y_j}$ (Lemma~\ref{lem:pauli} and $c_{y_j}=\ii Z_{<j}X_jZ_j$),
$(1\mp\lambda_jZ_j)c_{x_j}=c_{x_j}\mp\ii\lambda_jc_{y_j}$, which is the first identity; the
second is the same computation with $c_{y_j}$ and $Z_jc_{y_j}=-c_{y_j}Z_j=-\ii c_{x_j}$. The
rows of $\G$: $\Gamma(\rho_\lambda)=\D(\lambda)$ (Theorem~\ref{thm:mixture}) has the only
nonzero entries $\pm\lambda_j$ at $(x_j,y_j)$ and $(y_j,x_j)$.
\end{proof}
\begin{remark}\label{rem:kms-row}
For $|\lambda_j|<1$, solving the relation for $c_a\rho_\lambda$ gives the KMS-type form
$c_a\rho_\lambda=\rho_\lambda\sum_b\mathcal K_{ab}c_b$, where on $(x_j,y_j)$
\[
 \mathcal K=\begin{pmatrix}\cosh\varepsilon&-\ii\sinh\varepsilon\\ \ii\sinh\varepsilon&\cosh\varepsilon\end{pmatrix},
\] $\varepsilon=-2\artanh\lambda_j$, the block of $e^{-\ii h}$ for $h=\D(\varepsilon)$;
for $\lambda_j=-1$ it reads $2a_j\rho_\lambda=2\rho_\lambda a_j^*$, both sides zero because the
mode is empty, and for $\lambda_j=1$ the same with $a_j^*$. The single identity covers all
three cases; Theorem~\ref{thm:mixture-gaussian} needs no case distinction.
\end{remark}

\begin{theorem}[The product mixtures are Gaussian]\label{thm:mixture-gaussian}
For every $\lambda\in[-1,1]^n$ the product mixture $\rho_\lambda$ is Gaussian:
$\tr\bigl(\rho_\lambda\prod_i(v_i\cdot c)\bigr)=\W_{1-\ii\D(\lambda)}(v_1,\ldots,v_m)$ for all
$v_1,\ldots,v_m$.
\end{theorem}
\begin{proof}
We apply Lemma~\ref{lem:criterion} with $\rho=\rho_\lambda$ and $\G=\G(\rho_\lambda)=1-\ii\D(\lambda)$.
Condition (i) is Lemma~\ref{lem:parity}. For (ii) fix a mode $j$, a list $w_1,\ldots,w_k$ with
$k$ odd, $X=\prod_i(w_i\cdot c)$, and write $\langle Y\rangle=\tr(\rho_\lambda Y)$,
$A_a=\langle c_aX\rangle$, $B_a=\langle Xc_a\rangle$ and
$r_a=2\sum_i(-1)^{i-1}(w_i)_a\langle X_{\hat\imath}\rangle$ for $a\in\{x_j,y_j\}$.
Lemma~\ref{lem:move} (second identity, $u=e_a$; $k$ odd, so $(-1)^{k-i}=(-1)^{i-1}$) gives
$B_a=r_a-A_a$. Multiplying the first identity of Lemma~\ref{lem:mode} by $X$ on the right and
taking the trace, with $\tr(Y\rho_\lambda X)=\tr(\rho_\lambda XY)$,
\[
 B_{x_j}-\ii\lambda_jB_{y_j}=A_{x_j}+\ii\lambda_jA_{y_j}.
\]
Substituting $B_a=r_a-A_a$ for both labels, the terms $\pm\ii\lambda_jA_{y_j}$ cancel and
\[
 2A_{x_j}=r_{x_j}-\ii\lambda_jr_{y_j}
 =2\sum_i(-1)^{i-1}\bigl((w_i)_{x_j}-\ii\lambda_j(w_i)_{y_j}\bigr)\langle X_{\hat\imath}\rangle
 =2\sum_i(-1)^{i-1}(e_{x_j}^{\T}\G w_i)\langle X_{\hat\imath}\rangle
\]
by the row of $\G$ in Lemma~\ref{lem:mode}: this is (ii) for $a=x_j$. The second identity of
Lemma~\ref{lem:mode} gives in the same way $2A_{y_j}=r_{y_j}+\ii\lambda_jr_{x_j}$, which is (ii)
for $a=y_j$.
\end{proof}

\begin{corollary}[The Gaussian families]\label{cor:families}
\begin{enumerate}
\item[(a)] For $OO^{\T}=1$, $U$ implementing $O$ and $\lambda\in[-1,1]^n$, the state
$U^*\rho_\lambda U$ \vet{(as in the proof of Theorem~\ref{thm:realization}, which does not use $\det O=1$)} is Gaussian with
$\G=O(1-\ii\D(\lambda))O^{\T}$.
\item[(b)] Every Gibbs state $\rho_h$, $h^{\T}=-h$, and every thermal state
$\rho_{\tau h}$ is Gaussian.
\item[(c)] Every occupation state $|s\rangle\langle s|$ and every pure state
$U^*|s\rangle\langle s|U$ with $U$ implementing an orthogonal $O$ is Gaussian.
\end{enumerate}
\end{corollary}
\begin{proof}
(a) Theorems~\ref{thm:mixture-gaussian} and~\ref{thm:transport}. (b) Theorem~\ref{thm:thermal}:
$\rho_h=U^*\rho_\lambda U$ with $\lambda_j=-\tanh(\varepsilon_j/2)$ for a normal form
$h=O\D(\varepsilon)O^{\T}$ and $U$ implementing $O$; apply (a); $\rho_{\tau h}$ is the Gibbs
state of $\tau h$, and $(\tau h)^{\T}=-\tau h$. (c) $|s\rangle\langle s|=\rho_\lambda$ with $\lambda_j=\pm1$
(Theorem~\ref{thm:mixture}); apply (a).
\end{proof}

\subsection{Moments determine the state}
\label{sec:monomials}

\begin{definition}[Majorana monomials]\label{def:monomial}
Order the labels $x_1<\cdots<x_n<y_1<\cdots<y_n$. For a set $S$ of labels, $c_S$ is the product
of the $c_a$, $a\in S$, in increasing order; $c_\emptyset=1$.
\end{definition}

\begin{lemma}[Traces of monomials]\label{lem:monomial-trace}
\begin{enumerate}
\item[(a)] $\tr(c_S)=0$ for $S\neq\emptyset$.
\item[(b)] $c_S^*c_S=1$ and, for $S\neq T$, $c_S^*c_T=\pm c_{S\triangle T}$ with
$S\triangle T\neq\emptyset$ the symmetric difference. Hence $\tr(c_S^*c_T)=2^n[S=T]$.
\end{enumerate}
\end{lemma}
\begin{proof}
(a) If $|S|$ is odd, $\Pi_nc_S\Pi_n=-c_S$ by Lemma~\ref{lem:parity}, so
$\tr c_S=\tr(\Pi_nc_S\Pi_n)=-\tr c_S$. If $|S|\geq2$ is even, write $c_S=c_ac_{S'}$ with $a\in S$
and
$|S'|$ odd; the labels of $S'$ differ from $a$, so $c_{S'}c_a=(-1)^{|S'|}c_ac_{S'}=-c_ac_{S'}$
and $\tr(c_S)=\tr(c_{S'}c_a)=-\tr(c_S)$.
(b) Insertion rule: for a set $T$ of labels and a label $a$,
$c_ac_T=(-1)^{|\{t\in T:t<a\}|}c_{T\cup\{a\}}$ if $a\notin T$ and
$c_ac_T=(-1)^{|\{t\in T:t<a\}|}c_{T\setminus\{a\}}$ if $a\in T$: move $c_a$ to the right past
the factors $c_t$ with $t<a$ (one sign each, Theorem~\ref{thm:car}), where it either takes
its place in the order or meets $c_a$ and cancels by $c_ac_a=1$. Hence, by induction on $|S|$,
$c_Sc_T=\pm c_{S\triangle T}$ for all $S,T$: for $S=\emptyset$ this is $c_T=c_T$; otherwise
write $c_S=c_ac_{S'}$ with $a=\min S$, apply the
hypothesis to $c_{S'}c_T=\pm c_{S'\triangle T}$ and then the insertion rule to
$c_ac_{S'\triangle T}$, noting $S\triangle T=\{a\}\triangle(S'\triangle T)$. Since
$c_S^*=\pm c_S$ (reversing $|S|$ distinct anticommuting factors, $c_a^*=c_a$),
$c_S^*c_T=\pm c_{S\triangle T}$, with $c_S^*c_S=1$ exactly: $c_S^*c_S=c_{a_m}\cdots c_{a_1}c_{a_1}\cdots c_{a_m}$ for
$S=\{a_1<\cdots<a_m\}$, and the innermost pair $c_{a_1}c_{a_1}=1$ cancels, then the next, and
so on. The trace then follows from (a), $S\triangle T\neq\emptyset$ for $S\neq T$, and
$\tr1=2^n$.
\end{proof}

\begin{theorem}[The monomial basis]\label{thm:basis}
The $4^n$ monomials $c_S$ are linearly independent, hence a basis of the $4^n$-dimensional
space of matrices on $\mathcal F_n$, and every matrix $A$ satisfies
\[
 A=2^{-n}\sum_S\tr(Ac_S^*)\,c_S .
\]
In particular, if $\tr(Ac_S)=\tr(Bc_S)$ for all $S$, then $A=B$.
\end{theorem}
\begin{proof}
If $\sum_S\alpha_Sc_S=0$, multiply by $c_T^*$ on the left and take the trace:
$2^n\alpha_T=0$ by Lemma~\ref{lem:monomial-trace}(b). There are $2^{2n}=(2^n)^2$ subsets and
the matrices on the $2^n$-dimensional space $\mathcal F_n$ form a $(2^n)^2$-dimensional space, so the
independent family is a basis. Writing $A=\sum_S\alpha_Sc_S$ and taking
$\tr(Ac_T^*)=\sum_S\alpha_S\tr(c_Sc_T^*)=\sum_S\alpha_S\tr(c_T^*c_S)=2^n\alpha_T$ (cyclicity and
Lemma~\ref{lem:monomial-trace}(b)) gives the formula; the last statement follows since
$c_S^*=\pm c_S$.
\end{proof}

\begin{theorem}[A Gaussian state is determined by its covariance]\label{thm:unique}
If $\rho$ and $\sigma$ are Gaussian and $\Gamma(\rho)=\Gamma(\sigma)$, then $\rho=\sigma$.
\end{theorem}
\begin{proof}
For every $S=\{a_1<\cdots<a_m\}$, $\tr(\rho c_S)=\W_{\G(\rho)}(e_{a_1},\ldots,e_{a_m})$ by
Wick's rule, and $\G(\rho)=1-\ii\Gamma(\rho)=\G(\sigma)$; so $\tr(\rho c_S)=\tr(\sigma c_S)$ for
all $S$ and Theorem~\ref{thm:basis} gives $\rho=\sigma$.
\end{proof}
For an admissible $\Gamma$ write $\rho[\Gamma]$ for the Gaussian state with covariance
$\Gamma$, unique by Theorem~\ref{thm:unique}; it exists by Theorem~\ref{thm:realization} and
Corollary~\ref{cor:families}(a).

\begin{corollary}[Every Gaussian state is a rotated product mixture]\label{cor:converse}
Let $\rho$ be Gaussian with covariance $\Gamma=O\D(\eta)O^{\T}$, $OO^{\T}=1$, $\det O=1$,
$|\eta_j|\leq1$ (Theorem~\ref{thm:normal-form}), and let $U$ implement $O$. Then
\begin{enumerate}
\item[(a)] $\rho=U^*\rho_\eta U$;
\item[(b)] if $\Gamma$ is pure, $\rho=U^*|s\rangle\langle s|U$ with $s_j=[\eta_j=1]$, and
$\rho^2=\rho$;
\item[(c)] if $\Gamma$ is strictly admissible, $\rho=\rho_h$ for some $h^{\T}=-h$; and for
$h^{\T}=-h$, $h'^{\T}=-h'$ with $\Gamma(\rho_h)=\Gamma(\rho_{h'})$, $\rho_h=\rho_{h'}$.
\end{enumerate}
\end{corollary}
\begin{proof}
(a) $U^*\rho_\eta U$ is a density operator with covariance $\Gamma$
(Theorem~\ref{thm:realization}) and is Gaussian (Corollary~\ref{cor:families}(a)); apply
Theorem~\ref{thm:unique}. (b) For pure $\Gamma$, $|\eta_j|=1$ for all $j$
(Theorem~\ref{thm:normal-form}) and $\rho_\eta=|s\rangle\langle s|$ (Theorem~\ref{thm:mixture});
$\rho^2=U^*|s\rangle\langle s|UU^*|s\rangle\langle s|U=\rho$. (c) By
Theorem~\ref{thm:exp-realization} there is $h^{\T}=-h$ with $\Gamma(\rho_h)=\Gamma$; $\rho_h$ is
Gaussian (Corollary~\ref{cor:families}(b)); apply Theorem~\ref{thm:unique}, and likewise for
$\rho_h,\rho_{h'}$.
\end{proof}
\revf{For a pure covariance the Gaussian hypothesis of Theorem~\ref{thm:unique} can be dropped: two
density operators with the same pure covariance are equal (Corollary~\ref{cor:pure-determines}).}
